\begin{titlepage}
\centering
\vspace*{2.0cm}

{\Large\bfseries
Q-scan morphology of scattered-light transients\\[0.35em]
in LIGO O3 and O4a\par}

\vspace{0.9em}
{\large\itshape with an auxiliary-channel correlation analysis\par}

\vspace{2.4cm}
{\large Hardik Sher Singh\par}
\vspace{0.4em}
{\normalsize Department of Electronics and Communication Engineering\\
Thapar Institute of Engineering and Technology, Patiala, India\par}
\vspace{0.3em}
{\small\texttt{hsingh\_be23@thapar.edu}\par}

\vspace{2.8cm}
\begin{minipage}{0.85\textwidth}
\small
\noindent\textbf{Abstract.}
Scattered light is among the most persistent sources of transient noise in the
Advanced LIGO detectors. Stray light reflecting from a moving surface and
rejoining the main beam is phase modulated by $4\pi x/\lambda$, and when that
excursion exceeds a radian the noise appears in the time--frequency plane as a
train of harmonic \emph{arches} whose instantaneous frequency follows
$f_n = 2n|v|/\lambda$. Because the frequency depends on the magnitude of the
velocity, each cycle of the driving motion produces two arches, so the arch
repetition period is a direct measurement of the mechanical disturbance
responsible.

This work characterises the morphology of scattered-light transients in the
third and fourth LIGO observing runs using only public data. A labelled O3
sample is drawn from the Gravity Spy classification release; every trigger is
Q-transformed and reduced to a twenty-element feature vector describing the
arch train. The O3 sample is then tied to physical cause in two ways: by
rank-correlating the hourly trigger rate against seismometer band-limited RMS
in the earthquake, microseismic and anthropogenic bands, and by fitting the
$f_n = 2n|v|/\lambda$ predictor, trigger by trigger, to the Livingston OM3
suspension witness that the public O3 auxiliary release makes available at
\SI{256}{\hertz}.

Two properties of the public data determine the shape of the O4a half of the
analysis and are reported as findings in their own right: no Gravity Spy
classification catalogue has been released for O4a, and the O4a auxiliary
release contains no seismometer, accelerometer or suspension witness of any
kind. The O4a sample is therefore constructed from strain alone, using an
excess-power trigger finder and a classifier trained on the O3 features, and
the O3 comparison is made like-for-like by running the identical search on
both runs. The auxiliary analysis, being an O3 result, is what licenses
reading the O4a arch-period distribution as a statement about ground motion.

\vspace{0.6em}
\noindent\textbf{Keywords:} gravitational waves; LIGO; detector
characterisation; transient noise; scattered light; time--frequency analysis.
\end{minipage}

\vfill
{\small Analysis code and configuration:
\url{https://github.com/hardiksingh/LIGO_Glitch_Classification}\par}
\vspace{0.4em}
{\small \today\par}
\end{titlepage}
